% Section 7.7, from lectures/L07/appendix/b_exercises.md.

\section{Exercises}
\label{c7:sec:exercises}

\begin{aside}
\textbf{How to check your work.} Every exercise below that asks you to write a VHDL module comes
with a self-checking testbench under \file{lectures/L07/exercises/}. Write your module in its
directory \file{exercises/<module>/}, with the entity name and the \textbf{port order} the exercise
states, and then run it with GHDL \textemdash{} see \secref{c2:sec:testbenches}.

From this chapter onwards the designs reuse modules you wrote earlier, and they are not handed out a
second time: every exercise says which of your own files to copy in, and writes out the
\code{ghdl -a} line that needs them.
\end{aside}

% ------------------------------------------------------------------------------------------
\exerciseset{Timers by hand}

\exercise{A timer with a different target value}{Circuit}
\label{c7:ex:byhand}
Build a second timer by hand in CircuitVerse, following \secref{c7:sec:byhand}, but with the target
value \code{TICK_COUNT = 12} instead of the section's \code{10}.

\xpart{a} Retarget the comparator:
\begin{itemize}
  \item Write \code{12} in binary with 4 bits.
  \item The four XNOR gates stay exactly where they are. Say which of their constant inputs change,
    and to what.
  \item Say, in one sentence, why an equality comparator built from XNORs needs no gate added,
    removed or rewired in order to compare against a different value.
\end{itemize}

\xpart{b} Build it: reuse your 4-bit counter and comparator from the session, changing only the
constants. Set the clock period to \code{1000 ms} and confirm that \code{timeout} fires once every
thirteenth clock edge, for exactly one edge each time.

\xpart{c} Now remove the clear, so that \code{timeout} no longer forces the counter back to
\code{0000}. Predict, before you run it, how often \code{timeout} fires now. Run it and confirm.
Explain the result in terms of the counter's own wraparound (\secref{c6:sec:counter}).

\xpart{d} Break the comparator deliberately: disconnect one XNOR's output from the AND gate and hold
that AND input high instead. Predict which count values now raise \code{timeout}, before you
simulate, and confirm by stepping the counter through a whole cycle. This is what ``three of the
four bits match'' looks like, and it is the most common way for a hand-drawn comparator to be wrong.

\note[Hint]Counting \code{0} through \code{12} inclusive is thirteen edges, not twelve.
\Secref{c7:sec:concept}'s remark about \code{TICK_COUNT + 1} is the same off-by-one, and it is worth
sorting out here, where you can count the edges by eye, rather than later at 50~MHz.

\exercise{Four questions about the module}{Reasoning}
\label{c7:ex:questions}
Answer each of the questions below in one or two sentences, with reference to the \code{timer}
module in \secref{c7:sec:module}:

\xpart{a} Why is \code{timeout} a one-cycle pulse rather than a signal that stays high once the
target value has been reached?

\xpart{b} While a running timer's \code{enable} input is held low, what happens to its internal
counter from one clock edge to the next: is it cleared, does it pause, or does it carry on counting?

\xpart{c} Why does \code{timer} take \code{reset_s2_n} (the synchronised reset signal) rather than
the raw \code{reset_n} as its reset input?

\xpart{d} The CircuitVerse timer in Exercise~\ref{c7:ex:byhand} had to be rewired for the target
value to change, while the VHDL module needs no change at all. What replaces the rewiring, and when
is its value fixed?

% ------------------------------------------------------------------------------------------
\exerciseset{Timers in VHDL}

\exercise{Three tick counts}{Reasoning}
\label{c7:ex:ticks}
The DE0-CV board's system clock is \code{50 MHz}. For each target below, work out the
\code{TICK_COUNT} generic you would pass to the \code{timer} module:

\xpart{a} A \code{timeout} pulse once per second.

\xpart{b} A \code{timeout} pulse every \code{250} ms.

\xpart{c} A \code{timeout} pulse twenty times per second (\code{20 Hz}).

\note[Hint]The relationship between a period and its tick count is
$\mathit{TICK\_COUNT} = \mathit{seconds} \times 50\,000\,000 - 1$. The $-1$ term is the same
off-by-one Exercise~\ref{c7:ex:byhand} circled around: the counter runs \code{0} through
\code{TICK_COUNT} inclusive, so a period is \code{TICK_COUNT + 1} clock cycles, not
\code{TICK_COUNT}.

\exercise{\code{timer}}{VHDL}
\label{c7:ex:timer}
Write the \code{timer} module itself, as an entity named \code{timer}.

It is built live during the session and printed in full in \secref{c7:sec:module}, so this is no
puzzle. Write it anyway, without copying: every design from here to the end of the book instantiates
this module, both capstones included, and it is worth having built what you are about to reuse six
times.

The entity has one generic:

\begin{booktable}{@{}p{6em}p{4.4em}p{6em}L@{}}
  \thead{Generic} & \thead{Type} & \thead{Default} & \thead{Meaning} \\
  \midrule
  \code{TICK_COUNT} & \code{natural} & \code{50_000_000} & Number of ticks to count before a
    timeout. One second at the DE0-CV's \code{50 MHz} clock, rounded: by
    Exercise~\ref{c7:ex:ticks}'s rule the exact value is \code{49_999_999}, and 20~ns of error over
    a second is not worth the less readable constant. \\
\end{booktable}

and these ports:

\begin{booktable}{@{}p{5.6em}p{3.4em}p{5.2em}L@{}}
  \thead{Port} & \thead{Dir.} & \thead{Type} & \thead{Meaning} \\
  \midrule
  \code{clock} & in & \code{std_logic} & System clock. \\
  \code{reset_s2_n} & in & \code{std_logic} & Active-low, already synchronised reset. Asynchronous:
    it clears the counter and \code{timeout} with no clock edge involved. \\
  \code{enable} & in & \code{std_logic} & The timer's enable. While it is low the counter
    \textbf{holds} its value, and carries on from there rather than starting over. \\
  \code{timeout} & out & \code{std_logic} & One-cycle pulse, firing every \code{TICK_COUNT + 1}
    cycles while the timer is enabled. \\
\end{booktable}

The architecture should:
\begin{itemize}
  \item Declare the internal counter as \code{natural range 0 to TICK_COUNT}, so that the register
    is sized by the generic rather than left unconstrained.
  \item Use a single process, sensitive to \code{clock} and \code{reset_s2_n}, following the clocked
    process template from \secref{c3:sec:template}.
  \item Clear both the counter and \code{timeout} at reset, outside the clocked branch.
  \item Drive \code{timeout} low unconditionally at the top of the clocked branch, and high only at
    the edge that clears the counter. That pattern, a default followed by an override, is what makes
    the pulse exactly one cycle wide rather than two.
\end{itemize}

Then answer, in one sentence each:
\begin{itemize}
  \item \code{enable = '0'} must \textbf{pause} the counter, not clear it. Suppose you cleared it
    instead. Predict what the reference testbench reports, and which of the two properties above it
    is measuring when it does. Then run it and see. (\Chapref{c8:ch}'s \code{fsm_led} depends on the
    pause, which is why the testbench cares.)
  \item The counter counts \code{0} through \code{TICK_COUNT} inclusive, so a period is
    \code{TICK_COUNT + 1} cycles. Where in your code is that ``+ 1'' actually written? It does not
    appear as a literal anywhere.
  \item \Secref{c7:sec:byhand}'s hand-drawn \code{timeout} is a combinational decode of the count;
    yours is registered. Both have the same period. Which cycle does each fire on, and why does the
    difference matter to something downstream sampling \code{timeout}?
\end{itemize}

\modulefig[0.62]{lectures/L07/appendix/images/timer.png}

\begin{selfcheck}
\textbf{Self-check.} Name your entity \code{timer}, with the generic \code{TICK_COUNT}
(\code{natural}), the inputs \code{clock}, \code{reset_s2_n}, \code{enable} and the output
\code{timeout}, declared in that order; its testbench is in \file{exercises/timer/}. It runs three
full periods, so a timer firing only once does not pass, and it pauses mid-count to confirm that the
count continued rather than starting over.

\textbf{Keep this file.} Exercise~\ref{c7:ex:blinker} below reuses it, so does \code{walking_led} in
Exercise~\ref{c7:ex:walking}, and so do both of \chapref{c8:ch}'s capstones.
\end{selfcheck}

\exercise{\code{blinker}}{VHDL}
\label{c7:ex:blinker}
Write an entity named \code{blinker} that blinks an LED on and off once per second, by reusing the
\code{timer} module you just wrote rather than counting clock cycles yourself.

The entity has one generic:

\begin{booktable}{@{}p{6em}p{4.4em}p{6em}L@{}}
  \thead{Generic} & \thead{Type} & \thead{Default} & \thead{Meaning} \\
  \midrule
  \code{TICK_COUNT} & \code{natural} & \code{25_000_000} & Half a second at the DE0-CV's
    \code{50 MHz} clock. Expose it as a generic instead of baking the count in, so that a testbench
    can shorten the blink; the FPGA build only uses the default. \\
\end{booktable}

and these ports:

\begin{booktable}{@{}p{5em}p{3.4em}p{5.2em}L@{}}
  \thead{Port} & \thead{Dir.} & \thead{Type} & \thead{Meaning} \\
  \midrule
  \code{clock} & in & \code{std_logic} & System clock. \\
  \code{reset_n} & in & \code{std_logic} & Active-low reset, raw from the outside world.
    Asynchronous: activate it and \code{led} goes out immediately, without waiting for a clock
    edge. \\
  \code{led} & out & \code{std_logic} & Cleared to \code{'0'} while the reset signal is active, and
    left there until the first timeout after the reset signal is released. \\
\end{booktable}

Note that this entity takes the \textbf{raw} \code{reset_n} rather than an already synchronised
\code{reset_s2_n}. Synchronising it is part of the exercise, and it is what \secref{c4:sec:resetsync}
means by ``exactly one module in a design is the exception'': here that module is \code{blinker}
itself.

The architecture should:
\begin{itemize}
  \item Instantiate your \code{reset_sync} from Exercise~\ref{c4:ex:resetsync}, and use its output
    \code{reset_s2_n} as the reset for everything inside.
  \item Instantiate the \code{timer} module from \secref{c7:sec:module}, passing your own generic
    \code{TICK_COUNT} straight through to the timer's \code{TICK_COUNT}, and holding its
    \code{enable} input high.
  \item Toggle \code{led} in a synchronous process every time the timer's \code{timeout} pulse
    fires.
\end{itemize}

Do not reimplement the counting logic, or the synchroniser, yourself. Reuse \code{timer} and
\code{reset_sync} unchanged, via a \code{generic map} and a \code{port map}, following the
composition pattern in \secref{c7:sec:compose}. Explain why a timer period of \textbf{half} a second
gives a blink cycle of \textbf{one} second (a full period of on and then off).

Then answer, in one sentence each:
\begin{itemize}
  \item The LED has to be off during reset, and the testbench checks that it goes out with no clock
    edge in between. Which half of the assert-asynchronously\--release-synchronously pattern does
    that check, and which half does it not check?
  \item Your \code{timer} is reset by \code{reset_s2_n}, two edges after \code{reset_n} is released.
    What would go wrong if you fed it the raw \code{reset_n} instead? Name the failure, not just the
    rule.
\end{itemize}

\modulefig[0.56]{lectures/L07/appendix/images/blinker.png}

\begin{selfcheck}
\textbf{Self-check.} Name your entity \code{blinker}, with the generic \code{TICK_COUNT}
(\code{natural}), the inputs \code{clock}, \code{reset_n} and the output \code{led}, declared in that
order; its testbench is in \file{exercises/blinker/}. \textbf{Copy your own \file{reset_sync.vhd}
and \file{timer.vhd} into that directory} before you build; neither is handed out there, since you
wrote both:

\begin{shell}
cd lectures/L07/exercises/blinker
cp ../../../L04/exercises/reset_sync/reset_sync.vhd .   # the two you wrote yourself
cp ../timer/timer.vhd .
ghdl -a --std=93 reset_sync.vhd timer.vhd blinker.vhd blinker_tb.vhd
\end{shell}

The testbench overrides \code{TICK_COUNT} with a small value and checks that the LED toggles over
and over, stays put for exactly one timer period between toggles, and that the first toggle is late
by your synchroniser's two edges.
\end{selfcheck}

% ------------------------------------------------------------------------------------------
\exerciseset{The walking LED}

A single lit LED walking along a row, one position per timer tick, started and stopped with a
button. It is built live during the session and described in \secref{c7:sec:walking}.

This is the payoff from the last four chapters, and it writes almost no logic of its own: three
modules, two of them unchanged from \chapref{c4:ch} and one written here, composed.
Exercise~\ref{c7:ex:walking} builds it, and Exercise~\ref{c7:ex:variants} and
\ref{c7:ex:direction} change what you built.

\exercise{\code{walking_led}}{VHDL}
\label{c7:ex:walking}
Write an entity named \code{walking_led}.

The entity has two generics:

\begin{booktable}{@{}p{6em}p{4.4em}p{6em}L@{}}
  \thead{Generic} & \thead{Type} & \thead{Default} & \thead{Meaning} \\
  \midrule
  \code{LED_COUNT} & \code{natural} & \code{8} & How many LEDs the bit walks along. The row is this
    wide, and the walk wraps after this many steps. \\
  \code{TICK_COUNT} & \code{natural} & \code{25_000_000} & Passed straight through to your
    \code{timer}, so that the bit moves one position per timer period: one step per half second at
    the DE0-CV's \code{50 MHz} clock. \\
\end{booktable}

and these ports:

\begin{booktable}{@{}p{5em}p{3.4em}p{11.4em}L@{}}
  \thead{Port} & \thead{Dir.} & \thead{Type} & \thead{Meaning} \\
  \midrule
  \code{clock} & in & \code{std_logic} & System clock. \\
  \code{reset_n} & in & \code{std_logic} & Active-low reset, raw from the outside world.
    Synchronising it is your job, exactly as it was in \code{blinker}. \\
  \code{button_n} & in & \code{std_logic} & Active-low push button, raw. Every press starts or stops
    the walk. \\
  \code{led} & out & \code{std_logic_vector(LED_COUNT-1 downto 0)} & The LED row. Exactly one bit is
    lit at a time; at reset that is bit \code{0}. \\
\end{booktable}

The architecture composes three modules you have already written, and adds two processes:
\begin{itemize}
  \item Instantiate your \textbf{\code{reset_sync}}, and use its \code{reset_s2_n} as the reset for
    everything else in the design, the two instances below included.
  \item Instantiate your \textbf{\code{button_sync}} with \code{generic map(1)}, for a one-cycle
    pulse per press. Its button ports are vectors of \code{COUNT} bits while \code{button_n} here is
    a scalar, so bridge the width with a one-element \code{std_logic_vector(0 downto 0)} in each
    direction. \Secref{c7:sec:walking} shows it.
  \item Instantiate your \textbf{\code{timer}} with your own \code{TICK_COUNT}, drive its
    \code{enable} from the shift flag below and take its \code{timeout} as the shift tick.
  \item A process toggling \textbf{\code{shift_enable}} at every pulse from the button, cleared to
    \code{'0'} at reset. The walk therefore starts stopped.
  \item A process holding the \textbf{shift register} that drives \code{led}: at reset, load a single
    lit bit at position \code{0} and clear the rest; at every \code{timeout} from the timer, rotate
    one step towards the MSB and let the top bit wrap back into bit \code{0}. That is a PISO register
    (\secref{c6:sec:piso}) with its own serial output wired back to its serial input, which is what
    turns a one-shot shift-out into an endless walk.
\end{itemize}

Then answer, in one sentence each:
\begin{itemize}
  \item \code{shift_enable} drives the timer's \code{enable}, so stopping the walk stops the timer
    too. Exercise~\ref{c7:ex:timer} had you show that \code{enable} \textbf{pauses} the count rather
    than clearing it. What does an observer at the board see, at the press that restarts the walk,
    if your timer clears instead?
  \item \code{button_sync} is instantiated with \code{COUNT = 1} here and with \code{COUNT = 2} in
    \chapref{c8:ch}'s \code{fsm_led}, from the same file with no change. What would you have had to
    write instead if the width were fixed in the module rather than passed in as a generic?
\end{itemize}

\modulefig[0.7]{lectures/L07/appendix/images/walking_led.png}

\begin{selfcheck}
\textbf{Self-check.} Name your entity \code{walking_led}, with the generics \code{LED_COUNT} and
\code{TICK_COUNT} (both \code{natural}), the inputs \code{clock}, \code{reset_n}, \code{button_n} and
the output \code{led}, declared in that order; its testbench is in \file{exercises/walking_led/}. It
overrides both generics (\code{LED_COUNT = 4}, \code{TICK_COUNT = 3}) so that a full round takes a
handful of clock cycles rather than two seconds, and it checks the reset pattern, that nothing moves
before the first press, and that the number of shifts over a fixed window matches the rate the timer
sets. Copy in the three modules it composes first; the command is in \secref{c7:sec:walking}.
\end{selfcheck}

\exercise{Changing rate and width}{Simulation}
\label{c7:ex:variants}
Make the following changes one at a time, and rerun the testbench after each. Parts \textbf{a)} and
\textbf{b)} change only the testbench's constants \textemdash{} that they require no change at all in
your \file{walking_led.vhd} is the whole point of them. Part \textbf{c)} is the one that touches the
module.

\xpart{a} Change the walking speed:
\begin{itemize}
  \item Change the constant \code{TICK_COUNT} in the \textbf{testbench} from \code{3} to \code{6}.
  \item It should still pass, this time because the testbench adapts with you: its observation
    window is a fixed \code{PULSES * 6} cycles, and it derives the number of shifts it expects from
    \code{TICK_COUNT}, one shift per \code{TICK_COUNT + 1} cycles.
  \item Now find where that adaptation stops meaning anything. Keep raising \code{TICK_COUNT} past
    the length of the window and work out what the expected number becomes, and what the check then
    actually proves. A test that passes for the wrong reason is worth being able to recognise.
  \item Then explain why a changed tick count cannot break the \emph{design}, only slow it down.
\end{itemize}

\xpart{b} Change the width: change the testbench's constant \code{LED_COUNT} from \code{4} to
\code{8}. It should still pass, with no change at all in your \file{walking_led.vhd}. Explain what
the module would have had to look like for this to have required a change.

\xpart{c} Load two lit bits on opposite sides of the register at reset (e.g. \code{"10001000"} for
\code{LED_COUNT = 8}), and let them walk in step:
\begin{itemize}
  \item Point out the single line in your \file{walking_led.vhd} you need to change.
  \item \textbf{Predict} what the reference testbench will report, before you run anything.
  \item Then run it and see. Explain the result: is the design wrong, or is the testbench checking
    something that is no longer the specification?
\end{itemize}

\note[Hint]Both \code{LED_COUNT} and \code{TICK_COUNT} are generics, so parts \textbf{a)} and
\textbf{b)} require no change at all in the module's body. That is the whole point of a generic: one
module, many sizes, no changes.

\exercise{A second button that reverses the direction}{VHDL}
\label{c7:ex:direction}
The current design walks in one direction only (towards the MSB, where the top bit wraps back to the
bottom). Change your \file{walking_led.vhd} so that a \textbf{second} button reverses the walking
direction.

\xpart{a} Add a second, synchronised button input: run it through the same \code{button_sync}
instance, and pass \code{2} to its generic \code{COUNT} (\code{generic map(2)}, positionally, as
everywhere else in this book).

\xpart{b} Add a signal \code{direction}, and toggle it at the new button's edge, exactly as
\code{shift_enable} is toggled by the first button.

\xpart{c} Make the rotation in \code{SHIFT_PROCESS} depend on \code{direction}. When
\code{direction = '0'}, rotate towards the MSB (the book's convention from \secref{c6:sec:shift},
and what you wrote in Exercise~\ref{c7:ex:walking}):

\begin{vhdlcode}
shift_reg <= shift_reg(LED_COUNT-2 downto 0) & shift_reg(LED_COUNT-1);
\end{vhdlcode}

When \code{direction = '1'}, rotate the other way, towards bit 0:

\begin{vhdlcode}
shift_reg <= shift_reg(0) & shift_reg(LED_COUNT-1 downto 1);
\end{vhdlcode}

This is the one place in the book where a shift register deliberately goes against the convention:
reversing the direction is the whole point of the exercise. Note that the second expression is
exactly the mirror image \secref{c6:sec:shift} warns you to watch for in other people's code.

\xpart{d} Your entity's ports have now changed: \code{button_n} is a two-bit vector rather than a
single bit.
\begin{itemize}
  \item Explain why \file{walking_led_tb.vhd} can no longer be run against your version, and exactly
    what GHDL will complain about when you try.
  \item This is the ordinary cost of changing an interface, and it is worth feeling once: a testbench
    is written against a particular entity, and widening a port breaks every instantiation of it,
    including the ones you did not write.
  \item Describe, in prose, the sequence of presses you would use to convince yourself that the
    design works, and what you would expect the LEDs to do at each step. That sequence is exactly
    what a testbench for the new entity would have to drive.
\end{itemize}
